% Fichier engendré par tools/generate-tex.py à partir de 04-limites-continuite-derivation.
% Les démonstrations en français sont transcrites du fichier Lean (skill transcribe-lean-proof).
\documentclass[11pt,a4paper]{article}

% Compile aussi bien avec pdflatex qu'avec xelatex ou tectonic.
\usepackage{iftex}
\ifPDFTeX
  \usepackage[T1]{fontenc}
  \usepackage[utf8]{inputenc}
\else
  \usepackage{fontspec}
\fi
\usepackage[french]{babel}
\usepackage{amsmath,amssymb,amsthm}
\usepackage[margin=2.5cm]{geometry}
\usepackage[hidelinks]{hyperref}

\theoremstyle{definition}
\newtheorem{definition}{Définition}
\theoremstyle{plain}
\newtheorem{theoreme}{Théorème}
\newtheorem{lemme}[theoreme]{Lemme}

% Renvoi à la déclaration Lean, une ligne après la démonstration, aligné à droite.
\newcommand{\source}[2]{\par\smallskip\noindent\hfill{\footnotesize\href{#1}{\texttt{#2}}}\par\smallskip}

\title{Limites, continuité, dérivation}
\date{}

\begin{document}
\maketitle
% Les identifiants Lean en machine à écrire ne se coupent pas : on tolère des
% espacements plus lâches plutôt que des lignes qui débordent.
\sloppy

\section{LimitesContinuiteDerivation}

\noindent
Lycée \textemdash{} section \og{} Limites, continuité, dérivation \fg{}, niveau première
et terminale S. Les limites sont celles de Mathlib : \texttt{Tendsto f (nhds a) (nhds l)}
pour la limite en un point, \texttt{Tendsto f atTop (nhds l)} pour la limite en
\(+\infty\). Les trois premiers énoncés vérifient que ces écritures sont bien les
définitions \og{} \(\varepsilon\)\textendash{}\(\delta\) \fg{} et
\og{} \(\varepsilon\)\textendash{}A \fg{} du programme. La dérivabilité s'écrit
\texttt{HasDerivAt f d a} : \og{} \texttt{f} admet le nombre dérivé \texttt{d} en
\texttt{a} \fg{}.

\medskip\noindent
Le programme du lycée \emph{admet} l'essentiel de ce chapitre : opérations sur les
limites, théorème des valeurs intermédiaires, lien entre le signe de la dérivée et le
sens de variation. Ces résultats sont ici empruntés à Mathlib, où ils sont démontrés. La
démonstration Lean se réduit alors à citer le résultat de la bibliothèque ; les
transcriptions ci-dessous disent ce que ce résultat affirme et par quel argument il
s'obtient, en signalant à chaque fois que le travail est fait ailleurs.

\subsection{Définitions}

\noindent
Toutes les définitions du fichier sont posées ici, avant les démonstrations.

\begin{definition}
La \emph{tangente} à la courbe de $f$ au point d'abscisse $a$, lorsque le nombre dérivé
en $a$ vaut $d$, est la fonction affine
\[ T(x) = f(a) + d\,(x - a) . \]
\end{definition}
\source{https://github.com/Commutator-IO/learn-lean/blob/main/courses/02-lycee/04-limites-continuite-derivation/LimitesContinuiteDerivation.lean\#L21}{LimitesContinuiteDerivation.lean\#L21}

\subsection{Limite en un point, limite en l'infini}

\begin{theoreme}
La limite en un point est la définition du programme : $f$ tend vers $\ell$ en $a$ si et
seulement si
\[ \forall \varepsilon > 0,\ \exists \delta > 0,\ \forall x,\quad
   |x - a| < \delta \implies |f(x) - \ell| < \varepsilon . \]
\end{theoreme}

\begin{proof}
Mathlib définit la limite par les voisinages, mais donne pour tout espace métrique la
caractérisation équivalente : $f$ tend vers $\ell$ en $a$ si et seulement si, pour tout
$\varepsilon > 0$, il existe $\delta > 0$ tel que $\mathrm{dist}(x, a) < \delta$ entraîne
$\mathrm{dist}(f(x), \ell) < \varepsilon$. Sur $\mathbb{R}$ la distance est la valeur
absolue de la différence,
\[ \mathrm{dist}(x, y) = |x - y| , \]
de sorte que les deux énoncés ne diffèrent que par la notation. Dans chaque sens on prend
donc le même $\delta$ que celui fourni par l'autre membre, et on récrit les distances en
valeurs absolues.
\end{proof}
\source{https://github.com/Commutator-IO/learn-lean/blob/main/courses/02-lycee/04-limites-continuite-derivation/LimitesContinuiteDerivation.lean\#L28}{LimitesContinuiteDerivation.lean\#L28}

\begin{theoreme}
La limite en $+\infty$ est la définition avec $A$ : $f$ tend vers $\ell$ en $+\infty$ si
et seulement si
\[ \forall \varepsilon > 0,\ \exists A \in \mathbb{R},\ \forall x \ge A,\quad
   |f(x) - \ell| < \varepsilon . \]
\end{theoreme}

\begin{proof}
Même situation qu'au point précédent : la caractérisation métrique de la convergence le
long du filtre des grandes valeurs s'écrit avec $\mathrm{dist}(f(x), \ell) < \varepsilon$
pour $x \ge A$. En remplaçant la distance par $|f(x) - \ell|$, les deux membres sont le
même énoncé, mot pour mot.
\end{proof}
\source{https://github.com/Commutator-IO/learn-lean/blob/main/courses/02-lycee/04-limites-continuite-derivation/LimitesContinuiteDerivation.lean\#L43}{LimitesContinuiteDerivation.lean\#L43}

\begin{theoreme}
Asymptote horizontale : si $f$ tend vers $\ell$ en $+\infty$, l'écart entre la courbe et
la droite d'équation $y = \ell$ tend vers zéro,
\[ f(x) - \ell \xrightarrow[x \to +\infty]{} 0 . \]
\end{theoreme}

\begin{proof}
La fonction $x \mapsto f(x) - \ell$ est la différence de $f$, qui tend vers $\ell$, et de
la fonction constante $\ell$. Sa limite est donc $\ell - \ell = 0$.
\end{proof}
\source{https://github.com/Commutator-IO/learn-lean/blob/main/courses/02-lycee/04-limites-continuite-derivation/LimitesContinuiteDerivation.lean\#L50}{LimitesContinuiteDerivation.lean\#L50}

\begin{theoreme}
Asymptote verticale : $1/x \to +\infty$ lorsque $x \to 0$ par valeurs strictement
positives ; l'axe des ordonnées est asymptote à l'hyperbole de la fonction inverse.
\end{theoreme}

\begin{proof}
C'est un résultat de la bibliothèque sur les corps ordonnés : l'inverse tend vers
$+\infty$ à droite de zéro. L'argument est immédiat : pour $A > 0$ donné, tout $x$ tel que
$0 < x < 1/A$ vérifie $1/x > A$. Comme $1/x$ et $x^{-1}$ sont deux écritures de la même
fonction, l'énoncé s'y ramène.
\end{proof}
\source{https://github.com/Commutator-IO/learn-lean/blob/main/courses/02-lycee/04-limites-continuite-derivation/LimitesContinuiteDerivation.lean\#L56}{LimitesContinuiteDerivation.lean\#L56}

\begin{theoreme}
Asymptote oblique : l'écart entre la courbe de $x \mapsto ax + b + 1/x$ et la droite
d'équation $y = ax + b$ tend vers zéro en $+\infty$ :
\[ \bigl(ax + b + \tfrac{1}{x}\bigr) - (ax + b) \xrightarrow[x \to +\infty]{} 0 . \]
\end{theoreme}

\begin{proof}
La différence se simplifie :
\[ \bigl(ax + b + \tfrac{1}{x}\bigr) - (ax + b) = \frac{1}{x} , \]
et l'inverse tend vers zéro en $+\infty$ : pour $\varepsilon > 0$ donné, tout $x > 1/\varepsilon$
vérifie $0 < 1/x < \varepsilon$. Les coefficients $a$ et $b$ ont disparu : c'est bien
l'écart à la droite, et non la fonction elle-même, qui tend vers zéro.
\end{proof}
\source{https://github.com/Commutator-IO/learn-lean/blob/main/courses/02-lycee/04-limites-continuite-derivation/LimitesContinuiteDerivation.lean\#L62}{LimitesContinuiteDerivation.lean\#L62}

\subsection{Opérations sur les limites}

\begin{theoreme}
La limite d'une somme est la somme des limites : si $f \to \ell$ et $g \to \ell'$ en $a$,
alors $f + g \to \ell + \ell'$.
\end{theoreme}

\begin{proof}
Résultat de la bibliothèque, que le programme admet. Si $|f(x) - \ell| < \varepsilon/2$ et
$|g(x) - \ell'| < \varepsilon/2$ pour $x$ assez proche de $a$, alors l'inégalité
triangulaire donne
\[ \bigl|(f(x) + g(x)) - (\ell + \ell')\bigr|
   \le |f(x) - \ell| + |g(x) - \ell'| < \varepsilon . \]
\end{proof}
\source{https://github.com/Commutator-IO/learn-lean/blob/main/courses/02-lycee/04-limites-continuite-derivation/LimitesContinuiteDerivation.lean\#L69}{LimitesContinuiteDerivation.lean\#L69}

\begin{theoreme}
Celle d'un produit est le produit des limites : $f \times g \to \ell\ell'$.
\end{theoreme}

\begin{proof}
Même principe, en majorant l'écart par
\[ |f(x)g(x) - \ell\ell'| \le |f(x)|\,|g(x) - \ell'| + |\ell'|\,|f(x) - \ell| , \]
le premier facteur étant borné au voisinage de $a$ puisque $f$ y converge.
\end{proof}
\source{https://github.com/Commutator-IO/learn-lean/blob/main/courses/02-lycee/04-limites-continuite-derivation/LimitesContinuiteDerivation.lean\#L74}{LimitesContinuiteDerivation.lean\#L74}

\begin{theoreme}
Celle d'un quotient est le quotient des limites, $f/g \to \ell/\ell'$, tant que la limite
du dénominateur n'est pas nulle : $\ell' \ne 0$.
\end{theoreme}

\begin{proof}
L'hypothèse $\ell' \ne 0$ garantit que $g$ ne s'annule pas au voisinage de $a$, ce qui
permet de majorer l'écart des quotients. C'est exactement là que naissent les formes
indéterminées : sans cette hypothèse, l'énoncé est faux. En Lean, $\ell'$ nul n'empêcherait
pas d'écrire $\ell/\ell'$ \textemdash{} la bibliothèque pose $x/0 = 0$ pour que la division
soit une fonction totale \textemdash{} mais cette valeur de remplissage n'a aucun sens
mathématique, et l'hypothèse est bien nécessaire.
\end{proof}
\source{https://github.com/Commutator-IO/learn-lean/blob/main/courses/02-lycee/04-limites-continuite-derivation/LimitesContinuiteDerivation.lean\#L80}{LimitesContinuiteDerivation.lean\#L80}

\begin{theoreme}
Limite d'une composée : si $f \to b$ en $a$ et $g \to c$ en $b$, alors $g \circ f \to c$
en $a$.
\end{theoreme}

\begin{proof}
Les limites s'enchaînent : $f$ envoie les points proches de $a$ sur des points proches de
$b$, et $g$ envoie ceux-ci sur des points proches de $c$. En Lean, la composition des
limites est la transitivité de la relation \og{} envoyer un filtre dans un autre \fg{}.
\end{proof}
\source{https://github.com/Commutator-IO/learn-lean/blob/main/courses/02-lycee/04-limites-continuite-derivation/LimitesContinuiteDerivation.lean\#L86}{LimitesContinuiteDerivation.lean\#L86}

\subsection{Croissances comparées}

\begin{theoreme}
L'exponentielle l'emporte sur toute puissance : pour tout entier $n$,
\[ \frac{e^x}{x^n} \xrightarrow[x \to +\infty]{} +\infty . \]
\end{theoreme}

\begin{proof}
Résultat de la bibliothèque, cité tel quel dans le fichier Lean. Il repose sur la
minoration de l'exponentielle par un terme de sa série : pour $x \ge 0$,
\[ e^x \ge \frac{x^{n+1}}{(n+1)!} ,
   \qquad\text{d'où}\qquad \frac{e^x}{x^n} \ge \frac{x}{(n+1)!} \]
dès que $x > 0$, et le membre de droite tend vers $+\infty$.
\end{proof}
\source{https://github.com/Commutator-IO/learn-lean/blob/main/courses/02-lycee/04-limites-continuite-derivation/LimitesContinuiteDerivation.lean\#L92}{LimitesContinuiteDerivation.lean\#L92}

\begin{theoreme}
Le logarithme est négligeable devant l'identité :
\[ \frac{\ln x}{x} \xrightarrow[x \to +\infty]{} 0 . \]
\end{theoreme}

\begin{proof}
La bibliothèque énonce que $\ln$ est négligeable devant la fonction identité en
$+\infty$, ce qui signifie précisément que le quotient des deux tend vers zéro. Comme la
fonction identité vaut $x$ en $x$, le quotient est bien $\ln x / x$.
\end{proof}
\source{https://github.com/Commutator-IO/learn-lean/blob/main/courses/02-lycee/04-limites-continuite-derivation/LimitesContinuiteDerivation.lean\#L97}{LimitesContinuiteDerivation.lean\#L97}

\begin{theoreme}
En zéro, le facteur $x$ l'emporte sur le logarithme :
\[ x \ln x \xrightarrow[\substack{x \to 0 \\ x > 0}]{} 0 . \]
\end{theoreme}

\begin{proof}
La bibliothèque donne, pour tout exposant réel $r > 0$,
\[ \ln x \cdot x^{r} \xrightarrow[\substack{x \to 0 \\ x > 0}]{} 0 . \]
On l'applique à $r = 1$ : la puissance réelle $x^{1}$ vaut $x$, et le produit se lit dans
l'autre ordre, ce qui donne $x \ln x$. La restriction aux $x > 0$ n'est pas un artifice de
formalisation : le logarithme n'est défini que là.
\end{proof}
\source{https://github.com/Commutator-IO/learn-lean/blob/main/courses/02-lycee/04-limites-continuite-derivation/LimitesContinuiteDerivation.lean\#L102}{LimitesContinuiteDerivation.lean\#L102}

\subsection{Continuité}

\begin{theoreme}
La continuité en un point : la limite y existe et vaut l'image du point,
\[ f(x) \xrightarrow[x \to a]{} f(a) . \]
\end{theoreme}

\begin{proof}
Les deux membres sont le même énoncé : c'est ainsi que Mathlib définit la continuité en un
point. La transcription n'a rien à démontrer, seulement à constater que la définition de la
bibliothèque est celle du programme.
\end{proof}
\source{https://github.com/Commutator-IO/learn-lean/blob/main/courses/02-lycee/04-limites-continuite-derivation/LimitesContinuiteDerivation.lean\#L111}{LimitesContinuiteDerivation.lean\#L111}

\begin{theoreme}
Toute fonction dérivable en un point y est continue.
\end{theoreme}

\begin{proof}
Si $f$ admet le nombre dérivé $d$ en $a$, on écrit pour $x \ne a$
\[ f(x) - f(a) = (x - a) \times \frac{f(x) - f(a)}{x - a} . \]
Le second facteur tend vers $d$, donc reste borné au voisinage de $a$, tandis que le
premier tend vers zéro : le produit tend vers zéro, c'est-à-dire $f(x) \to f(a)$.

\medskip\noindent
La réciproque est fausse \textemdash{} la valeur absolue est continue en zéro sans y être
dérivable, ses taux d'accroissement valant $-1$ à gauche et $+1$ à droite. Ce
contre-exemple est mentionné dans le fichier Lean mais n'y est pas formalisé.
\end{proof}
\source{https://github.com/Commutator-IO/learn-lean/blob/main/courses/02-lycee/04-limites-continuite-derivation/LimitesContinuiteDerivation.lean\#L116}{LimitesContinuiteDerivation.lean\#L116}

\subsection{Théorème des valeurs intermédiaires}

\begin{theoreme}
Théorème des valeurs intermédiaires : si $f$ est continue sur $[a\,;b]$ avec $a \le b$,
alors pour tout $k$ compris entre $f(a)$ et $f(b)$ il existe $c \in [a\,;b]$ tel que
$f(c) = k$.
\end{theoreme}

\begin{proof}
Le programme admet ce théorème ; Mathlib le démontre sous la forme d'une inclusion
d'ensembles,
\[ [f(a)\,;f(b)] \subseteq f\bigl([a\,;b]\bigr) , \]
c'est-à-dire que toute valeur intermédiaire est atteinte. Dire que $k$ appartient à
l'image, c'est exactement dire qu'il existe $c$ dans $[a\,;b]$ avec $f(c) = k$ : la
démonstration consiste à lire l'appartenance à l'image sous cette forme.

\medskip\noindent
Le contenu du théorème est la complétude de $\mathbb{R}$, jamais énoncée au lycée : c'est
elle qui interdit à la courbe de \og{} sauter \fg{} par-dessus la valeur $k$. Mathlib
l'obtient de la connexité de l'intervalle.
\end{proof}
\source{https://github.com/Commutator-IO/learn-lean/blob/main/courses/02-lycee/04-limites-continuite-derivation/LimitesContinuiteDerivation.lean\#L123}{LimitesContinuiteDerivation.lean\#L123}

\begin{theoreme}
Corollaire : si de plus $f$ est strictement croissante sur $[a\,;b]$, la solution de
$f(x) = k$ y est unique.
\end{theoreme}

\begin{proof}
L'existence est le théorème précédent : il fournit $c \in [a\,;b]$ tel que $f(c) = k$.
Pour l'unicité, soit $y \in [a\,;b]$ un autre point avec $f(y) = k$. Une fonction
strictement croissante est injective sur son intervalle : si $y \ne c$, l'un des deux est
strictement plus petit que l'autre, donc son image l'est aussi, et $f(y) \ne f(c)$. Comme
$f(y) = k = f(c)$, il reste $y = c$.
\end{proof}
\source{https://github.com/Commutator-IO/learn-lean/blob/main/courses/02-lycee/04-limites-continuite-derivation/LimitesContinuiteDerivation.lean\#L130}{LimitesContinuiteDerivation.lean\#L130}

\subsection{Nombre dérivé et tangente}

\begin{theoreme}
Le nombre dérivé est la limite du taux d'accroissement : $f$ admet le nombre dérivé $d$
en $a$ si et seulement si
\[ \frac{f(x) - f(a)}{x - a} \xrightarrow[\substack{x \to a \\ x \ne a}]{} d . \]
\end{theoreme}

\begin{proof}
Mathlib appelle \emph{pente} de $f$ entre $a$ et $x$ la quantité
\[ \mathrm{pente}(x) = (x - a)^{-1}\,\bigl(f(x) - f(a)\bigr) , \]
qui est le taux d'accroissement écrit avec un inverse plutôt qu'une barre de fraction :
les deux fonctions sont égales, point par point. La bibliothèque démontre que la
dérivabilité équivaut à la convergence de la pente lorsque $x$ tend vers $a$ en restant
distinct de $a$ ; il suffit donc de substituer une écriture à l'autre.

\medskip\noindent
La restriction $x \ne a$ est indispensable : en $x = a$ le taux d'accroissement s'écrit
$0/0$. Lean donne à ce quotient la valeur zéro pour rendre la division totale, mais cette
valeur ne veut rien dire ici, et le voisinage épointé l'écarte.
\end{proof}
\source{https://github.com/Commutator-IO/learn-lean/blob/main/courses/02-lycee/04-limites-continuite-derivation/LimitesContinuiteDerivation.lean\#L143}{LimitesContinuiteDerivation.lean\#L143}

\begin{theoreme}
Équation de la tangente : la droite d'équation $y = f(a) + d\,(x - a)$ passe par le point
de la courbe d'abscisse $a$, et a le nombre dérivé pour coefficient directeur.
\end{theoreme}

\begin{proof}
En $x = a$, le second terme s'annule et il reste $T(a) = f(a)$ : la droite passe bien par
le point de la courbe.

Pour le coefficient directeur, on dérive $x \mapsto f(a) + d\,(x - a)$ en un point
quelconque : la fonction $x \mapsto x - a$ a pour dérivée $1$, la multiplier par la
constante $d$ donne $d \times 1 = d$, et ajouter la constante $f(a)$ ne change rien. La
tangente est donc la fonction affine de coefficient directeur $d$ qui coïncide avec $f$ en
$a$.
\end{proof}
\source{https://github.com/Commutator-IO/learn-lean/blob/main/courses/02-lycee/04-limites-continuite-derivation/LimitesContinuiteDerivation.lean\#L153}{LimitesContinuiteDerivation.lean\#L153}

\subsection{Dérivées usuelles}

\begin{theoreme}
Dérivées des fonctions de référence :
\[ (x^n)' = n\,x^{n-1}, \qquad
   \Bigl(\frac{1}{x}\Bigr)' = -\frac{1}{x^{2}} \ (x \ne 0), \qquad
   (\sqrt{x}\,)' = \frac{1}{2\sqrt{x}} \ (x > 0), \]
\[ \sin' = \cos, \qquad \cos' = -\sin, \qquad \exp' = \exp, \qquad
   (\ln x)' = \frac{1}{x} \ (x \ne 0). \]
\end{theoreme}

\begin{proof}
Les sept dérivées sont celles de la bibliothèque, citées une à une ; deux d'entre elles
demandent une mise en forme.

\smallskip\noindent
\emph{L'inverse.} Mathlib donne $(x^{-1})' = -(x^{2})^{-1}$ pour $x \ne 0$. C'est la même
chose que $-1/x^{2}$, une fois les inverses écrits en fractions.

\smallskip\noindent
\emph{Le logarithme.} La bibliothèque donne $x^{-1}$, que l'on récrit $1/x$.

\smallskip\noindent
Deux remarques sur les hypothèses. La racine carrée de Lean est totale \textemdash{} elle
vaut zéro sur les négatifs \textemdash{} et sa formule de dérivation exige $x > 0$, comme
au lycée. Pour la puissance, l'exposant $n - 1$ est une soustraction d'entiers naturels,
donc tronquée : pour $n = 0$ elle vaut $0$ au lieu de $-1$, mais le facteur $n = 0$ annule
le produit et la formule reste juste, la dérivée d'une constante étant nulle.
\end{proof}
\source{https://github.com/Commutator-IO/learn-lean/blob/main/courses/02-lycee/04-limites-continuite-derivation/LimitesContinuiteDerivation.lean\#L165}{LimitesContinuiteDerivation.lean\#L165}

\subsection{Opérations sur les dérivées}

\begin{theoreme}
Si $u$ et $v$ ont pour nombres dérivés $u'$ et $v'$ en $x$, alors
\[ (u + v)' = u' + v', \qquad (uv)' = u'v(x) + u(x)v', \]
\[ \Bigl(\frac{1}{v}\Bigr)' = -\frac{v'}{v(x)^{2}}, \qquad
   \Bigl(\frac{u}{v}\Bigr)' = \frac{u'v(x) - u(x)v'}{v(x)^{2}}, \]
les deux dernières là où $v$ ne s'annule pas.
\end{theoreme}

\begin{proof}
La somme, le produit et le quotient sont les règles de la bibliothèque, citées
directement.

L'inverse s'en déduit sans nouvelle règle : c'est le quotient de la fonction constante
égale à $1$, de dérivée nulle, par $v$. La formule du quotient donne alors
\[ \frac{0 \times v(x) - 1 \times v'}{v(x)^{2}} = -\frac{v'}{v(x)^{2}} , \]
ce qui est bien la formule annoncée.
\end{proof}
\source{https://github.com/Commutator-IO/learn-lean/blob/main/courses/02-lycee/04-limites-continuite-derivation/LimitesContinuiteDerivation.lean\#L184}{LimitesContinuiteDerivation.lean\#L184}

\subsection{Dérivée d'une composée}

\begin{theoreme}
Dérivée d'une composée : si $u$ a pour nombre dérivé $u'$ en $x$ et $v$ pour nombre
dérivé $v'$ en $u(x)$, alors
\[ (v \circ u)'(x) = v'\bigl(u(x)\bigr) \times u'(x) . \]
\end{theoreme}

\begin{proof}
C'est la règle de dérivation en chaîne de la bibliothèque, citée telle quelle : le facteur
$u'$ mesure la vitesse à laquelle $u$ parcourt son image, et $v'$ la sensibilité de $v$ au
point atteint.
\end{proof}
\source{https://github.com/Commutator-IO/learn-lean/blob/main/courses/02-lycee/04-limites-continuite-derivation/LimitesContinuiteDerivation.lean\#L199}{LimitesContinuiteDerivation.lean\#L199}

\begin{theoreme}
Cas particulier $x \mapsto u(ax + b)$, dont la dérivée est $a\,u'(ax + b)$.
\end{theoreme}

\begin{proof}
La fonction affine $t \mapsto at + b$ a pour dérivée $a$ : l'identité a pour dérivée $1$,
la multiplier par $a$ donne $a \times 1 = a$, et la constante $b$ ne change rien. La règle
de la composée donne alors $u' \times a$ pour dérivée de $x \mapsto u(ax + b)$, que l'on
écrit $a\,u'$ en échangeant les facteurs.
\end{proof}
\source{https://github.com/Commutator-IO/learn-lean/blob/main/courses/02-lycee/04-limites-continuite-derivation/LimitesContinuiteDerivation.lean\#L203}{LimitesContinuiteDerivation.lean\#L203}

\subsection{Signe de la dérivée et sens de variation}

\begin{theoreme}
Si $f$ est continue sur $[a\,;b]$ et si $f' > 0$ sur $]a\,;b[$, alors $f$ est strictement
croissante sur $[a\,;b]$.
\end{theoreme}

\begin{proof}
Le programme admet ce résultat ; Mathlib le démontre pour une partie convexe quelconque de
$\mathbb{R}$, à partir du théorème des accroissements finis : pour $x < y$ dans
l'intervalle, il existe $c$ entre eux tel que
\[ f(y) - f(x) = (y - x)\,f'(c) > 0 . \]
Le segment $[a\,;b]$ est convexe et son intérieur est l'intervalle ouvert $]a\,;b[$ : les
hypothèses du fichier sont donc exactement celles du résultat cité, une fois l'intérieur
identifié.
\end{proof}
\source{https://github.com/Commutator-IO/learn-lean/blob/main/courses/02-lycee/04-limites-continuite-derivation/LimitesContinuiteDerivation.lean\#L215}{LimitesContinuiteDerivation.lean\#L215}

\begin{theoreme}
Si $f$ est continue sur $[a\,;b]$ et si $f' < 0$ sur $]a\,;b[$, alors $f$ est strictement
décroissante sur $[a\,;b]$.
\end{theoreme}

\begin{proof}
Même argument avec l'inégalité renversée : $f(y) - f(x) = (y - x)f'(c) < 0$ pour $x < y$.
\end{proof}
\source{https://github.com/Commutator-IO/learn-lean/blob/main/courses/02-lycee/04-limites-continuite-derivation/LimitesContinuiteDerivation.lean\#L222}{LimitesContinuiteDerivation.lean\#L222}

\subsection{Extremum local}

\begin{theoreme}
Si $f$ admet le nombre dérivé $d$ en $a$ et présente en $a$ un maximum local, alors
$d = 0$.
\end{theoreme}

\begin{proof}
C'est le lemme de Fermat, cité à la bibliothèque. Son argument : au voisinage de $a$ on a
$f(x) \le f(a)$, donc le taux d'accroissement
\[ \frac{f(x) - f(a)}{x - a} \]
est négatif ou nul pour $x > a$, positif ou nul pour $x < a$. Sa limite $d$ est donc à la
fois $\le 0$ et $\ge 0$, c'est-à-dire nulle. L'hypothèse que le point est intérieur est
portée par la notion de maximum \emph{local}, qui compare $f(a)$ aux valeurs prises de part
et d'autre.
\end{proof}
\source{https://github.com/Commutator-IO/learn-lean/blob/main/courses/02-lycee/04-limites-continuite-derivation/LimitesContinuiteDerivation.lean\#L232}{LimitesContinuiteDerivation.lean\#L232}

\begin{theoreme}
La réciproque est fausse : la fonction cube a un nombre dérivé nul en zéro,
\[ (x^{3})'\big|_{x=0} = 3 \times 0^{2} = 0 , \]
sans y présenter de maximum local.
\end{theoreme}

\begin{proof}
Le nombre dérivé se lit sur la formule $(x^n)' = n x^{n-1}$ prise en $n = 3$ et $x = 0$ :
il vaut $3 \times 0^{2} = 0$.

Supposons maintenant que zéro soit un maximum local : on aurait $x^{3} \le 0^{3} = 0$ pour
tout $x$ assez proche de zéro. Cette propriété resterait vraie en n'approchant zéro que par
la droite ; or tout voisinage de zéro à droite contient un $x > 0$, et un tel $x$ vérifie
$x^{3} > 0$. Les deux inégalités $x^3 \le 0$ et $x^3 > 0$ se contredisent.
\end{proof}
\source{https://github.com/Commutator-IO/learn-lean/blob/main/courses/02-lycee/04-limites-continuite-derivation/LimitesContinuiteDerivation.lean\#L237}{LimitesContinuiteDerivation.lean\#L237}

\subsection{Dérivée seconde, convexité, point d'inflexion}

\begin{theoreme}
Si $f$ est dérivable ainsi que $f'$ sur $[a\,;b]$, et si $f'' \ge 0$ sur $[a\,;b]$, alors
$f$ est convexe sur $[a\,;b]$.
\end{theoreme}

\begin{proof}
Résultat de la bibliothèque, cité pour la partie convexe $[a\,;b]$. Son argument : une
dérivée seconde positive rend $f'$ croissante, et une fonction dont la dérivée croît est
convexe \textemdash{} sa courbe reste au-dessous de ses cordes. Les hypothèses du fichier
sont celles du résultat cité, à ceci près qu'elles portent sur le segment fermé tout entier
plutôt que sur son intérieur.
\end{proof}
\source{https://github.com/Commutator-IO/learn-lean/blob/main/courses/02-lycee/04-limites-continuite-derivation/LimitesContinuiteDerivation.lean\#L254}{LimitesContinuiteDerivation.lean\#L254}

\begin{theoreme}
Point d'inflexion : la fonction cube a pour dérivée $3x^{2}$ et pour dérivée seconde
$6x$, qui s'annule en zéro en changeant de signe ; la fonction est convexe sur
$[0\,;+\infty[$.
\end{theoreme}

\begin{proof}
La première dérivée se lit sur la formule des puissances : $(x^{3})' = 3x^{3-1} = 3x^{2}$,
l'exposant $3 - 1$ se calculant dans les entiers naturels.

Pour la seconde, on dérive $t \mapsto 3t^{2}$ : la puissance donne $2x^{2-1} = 2x$, et le
facteur constant $3$ se conserve, d'où
\[ 3 \times (2x) = 6x . \]
Cette expression est négative avant zéro et positive après : la dérivée seconde change de
signe en zéro, ce qui est la définition d'un point d'inflexion.

Enfin, la convexité sur les positifs est celle des puissances entières, démontrée dans la
bibliothèque pour tout exposant.

\medskip\noindent
Le fichier ne formalise pas la concavité sur $]-\infty\,;0]$, qui achèverait le
changement de convexité ; seul le changement de signe de la dérivée seconde y est établi.
\end{proof}
\source{https://github.com/Commutator-IO/learn-lean/blob/main/courses/02-lycee/04-limites-continuite-derivation/LimitesContinuiteDerivation.lean\#L264}{LimitesContinuiteDerivation.lean\#L264}

\vfill
\noindent\rule{0.35\linewidth}{0.4pt}\par
\noindent{\footnotesize Leonardo de Moura et Sebastian Ullrich,
\emph{The Lean~4 Theorem Prover and Programming Language},
\emph{Automated Deduction — CADE~28}, Springer, 2021, p.~625--635,
\href{https://doi.org/10.1007/978-3-030-79876-5_37}{doi:10.1007/978-3-030-79876-5\_37}.\par}

\end{document}
